\documentclass[11pt]{article}
\usepackage{amsmath,hyperref}
\title{A named right can be used before an older operation drains}
\date{}
\begin{document}\maketitle

\paragraph{Result and scope.}
In an authored native warehouse case, both old worlds have the same
current cargo, carrier poses and last committed placement generation.
One live old operation will touch only $X$; another will exchange $X$
and $Y$. A true local label of $Y$ is not enough to start the same new
withdrawal in both. An enforceable disjoint predecessor scope, or a later
native receipt that closes the remaining old use of $Y$, permits useful
work before the entire old operation drains. The implemented
same-information direct policy ties. There is no deployment claim.

\paragraph{Actual common prefix and permissions.}
Start from the earlier native deliveries and placements. Execute old
$X$ pickup, transport to a temporary buffer and physical unload.
This prefix is identical for the two already accepted operation
scopes. $Y$ still sits in the original bay; no new placement
generation has committed. The new right names $Y$ and may be
exercised by a helper at a qualified eligible bay. It is subordinate
to preserving the old operation. The read and guard use the
current native identity, not its archived placement. Certificates of
predecessor scope and time validity are inputs to the fixture;
the actor is not given the old operation's hidden branch when
that qualification is missing, incomplete or expired.

The old re-seat returns $X$ from the buffer to its original bay.
Its prescribed route and shelf are disjoint from the new helper's
movements and use of $Y$. The old exchange must first lift $Y$
from its original bay, move it into the former $X$ bay and clear
its carrier; it then returns $X$ into the other bay. The code
actually checks the course of these native moves, identity,
buffer and unload effects. In an isolated negative control,
removing only the correctly named $Y$ under the new request
without checking the old scope leaves the old exchange's
subsequent declared source bay empty. The native old lift gets
no shelf. The control stops there; the runtime scope-aware policy
never admits this conflict. A reroute to an uncontrolled refuge
or automatic reauthorization of that old duty is not assumed.

\paragraph{Local admission and exact timing.}
The early new action is safe at the shared prefix when every
admitted old continuation is covered by an operative scope
that excludes the new action's effects. Here those effects
are the $Y$ withdrawal and its prescribed route; geometry and
controller checks supply route independence where claimed.
Different object names alone do not establish physical
noninterference. A missing, partial or expired scope cannot
exclude the exchange. Once the old exchange's first phase
really clears its carrier and leaves $Y$ in its revised
location, the old remaining authorized write set excludes
$Y$. In the phase-qualified variant a source receipt marks
that fact. The revised named snapshot is checked there.
The earlier whole-operation generation may still be unchanged.
The latest qualifying scope must be operative through the actual
new action, not merely at receipt time.

A staged new helper takes six native steps through approach,
pickup, turn, retreat and unload into its refuge. With complete
re-seat scope it can start at tick~1 and finish by tick~6 while
old $X$ completes at tick~16. With exchange, the old $Y$
carrier actually clears at tick~12, so the phase receipt
makes the revised local removal feasible by tick~18 even though
old $X$ finishes at tick~28. At the tick~20 deadline, waiting
for the whole old operation misses that new removal. With
a later horizon the whole-operation wait can also complete.
The exact direct baseline uses the earliest of the same
authorized cuts, rather than an omniscient clearance or
a disabled local actor. No earliest-time claim is made over
different motions, a broader authorization, or stepping into
a departing holder's cell before the available clearance proof.

Later delivery is a distinct native effect. Its route identity
is fixed at evaluation; relocating $Y$ in the exchange may
clear the new bay but not the original one. The due path
runs after the earlier relocation deadline and the declared
old drain, under its separate release, and actually receives
its native goal reward only when its own bay was cleared.
A different path is not substituted for the assigned one.
An independent carrier pays in every admitted fixture run.

\paragraph{Comparisons and boundary.}
This is physical evidence of the distinction underlying
\href{https://cs.nyu.edu/~shasha/papers/chopping.pdf}
{transaction chopping} and separate local phases in
\href{https://www.cs.princeton.edu/techreports/1987/070.pdf}
{Sagas}; their database machinery is not tested or claimed inferior.
The control that holds all new work until old drain is a
conservative negative comparison, not the strongest attainable
scheduler. The matched direct scoped controller reproduces all
new effects at the same cuts. Reversible source phases, compensation,
an unknown raw plant law and concurrent issuance by an actor
not covered by the scope would need different obligations.

The replay checks 384 authored source regimes and two explicit
incompatible-order controls. Three Lean results check scope
restriction, the incompatible-history obstruction and truly
disjoint product updates. They do not prove that a production
issuer can certify the future phase scope, or abstract native motion.
No independent case, fault disposition, R5/I1 acceptance or
statistical reliability is inferred.
\end{document}
